\section{Elipse de acción y reciprocidad modular}
\label{sec:ellipse}
Fijadas las salidas $A,C^*,\hbar>0$ en la coordenatización
$|C^*/A|<1$, definimos
\[
\eta_{\rm el}=\operatorname{artanh}(C^*/A),\quad
a_S=\sqrt{2\hbar}\,e^{\eta_{\rm el}/2},\quad
b_S=\sqrt{2\hbar}\,e^{-\eta_{\rm el}/2}.
\]
Las coordenadas canónicas equilibradas tienen unidad de raíz de acción;
no se suman aquí posición y momento sin una coordenatización dimensional.
Se cumplen
\[
a_Sb_S=2\hbar,\qquad
\tau_{\rm el}=i\,\frac{b_S}{a_S},\qquad
R_{\rm el}=\sqrt{\frac{b_S}{a_S}}
=\left(\frac{A-C^*}{A+C^*}\right)^{1/4}.
\]
La raíz positiva determina la composición $\tau_{\rm el}=iR_{\rm el}^2$,
y permite recuperar la razón angular:
\[
\frac{C^*}{A}=\frac{1-R_{\rm el}^4}{1+R_{\rm el}^4}.
\]
La inversión de forma envía $R_{\rm el}$ a $R_{\rm el}^{-1}$ y
$\tau_{\rm el}$ a $-1/\tau_{\rm el}$.

\begin{proposition}[Energía coincidente, acción distinguible]
Sea
\[
D(\eta)=\begin{pmatrix}e^\eta&0\\0&e^{-\eta}\end{pmatrix},
\quad
S_2=\begin{pmatrix}0&1\\1&0\end{pmatrix},
\quad
J_2=\begin{pmatrix}0&-1\\1&0\end{pmatrix}.
\]
Ambos mapas satisfacen $T^{\mathsf T}D(-\eta)T=D(\eta)$.
Sin embargo, para $\omega=\dd Q\wedge\dd P$,
\[
S_2^*\omega=-\omega,\qquad J_2^*\omega=\omega.
\]
Si $\mathcal I(\gamma)=\oint_\gamma P\,\dd Q$ sobre una curva cerrada,
\[
\mathcal I(S_2\gamma)=-\mathcal I(\gamma),\qquad
\mathcal I(J_2\gamma)=\mathcal I(\gamma).
\]
\end{proposition}
\begin{proof}
La conjugación intercambia las entradas diagonales de $D$.
Los determinantes de $S_2,J_2$ son $-1,1$, respectivamente, lo que
prueba la transformación de $\omega$. Para $\lambda=P\,\dd Q$,
\[
S_2^*\lambda=\dd(PQ)-\lambda,\qquad
J_2^*\lambda=\lambda-\dd(PQ).
\]
Integrar la diferencial exacta sobre la curva cerrada completa la prueba.
\end{proof}

La energía cuadrática y el módulo no bastan para distinguir ambos
transportes. La orientación de la acción sí los distingue, incluso
en la forma circular. El marcado geométrico es así parte de la lectura
recuperable, no un adorno añadido al valor.
